\chapter{Conclusion}
\label{chap:conclusion}

\section{Answers to the Research Questions}
\paragraph{\ref{rq1}} A qualified yes, for WebGen-LM-7B with Bolt.diy on WebGen-Bench. The final code memory raises accuracy from 28.5\% to 35.2\% (+7.1, CI excludes zero) and cuts zero-score tasks from 30 to 20, with gains of 7.4 and 5.5 points in two shuffled orders; most of the gain comes from sites that now render. \todo{Status: summarise all orders once seed 47 and interleaved finish.}

\paragraph{\ref{rq2}} The small generator followed short code fragments from its own earlier websites but rarely followed textual lessons; text memory fell behind No Memory late in the sequence (a cross-version, not controlled, comparison). The checklist alone added little (+3.4, interval includes zero). Fragments must be partial: a complete file was copied as a finished answer. \todo{Status: add the fixed-prefix result (queued after seed 47).}

\paragraph{\ref{rq3}} Code memory keeps its advantage over the sequence (35.3\% vs 29.5\% on tasks 31--101) while text memory declines, consistent with, but not proof of, learning from experience. Whole-file examples caused scaffolds when the order changed; with verified fragments the scaffolds disappeared and the gain was positive in both kinds of order, though robustness is not yet established. The whole-file design scored higher only in the original order (40.0\%). \todo{Seed 47 and interleaved order with the final method.}

\section{Future Work}
Two further experiments would address the main remaining questions:
\begin{itemize}
  \item \textbf{Held-out evaluation.} Because the method was developed while observing results on all 101 tasks, a clean test needs tasks that were never seen: a held-out domain, or a second website-generation benchmark, with the method frozen beforehand.
  \item \textbf{Controlled form-of-memory comparison.} Keep retrieval, trigger and token budget fixed and vary only the form: textual rule, pseudocode, short fragment, complete file, plus random-code and length-matched text controls.
\end{itemize}
Further directions are repeating each order more than once, to separate order effects from run-to-run variation, and testing larger generators of the same family, with a text-memory arm, to see whether the preference for code depends on model size.
